% Generated draft: requires mathematical and source review.
\documentclass[11pt]{article}
\usepackage[margin=25mm]{geometry}
\usepackage{fontspec}
\setmainfont{Latin Modern Roman}
\usepackage{amsmath,amssymb,mathtools,unicode-math}
\setmathfont{Latin Modern Math}
\usepackage{enumitem,needspace,xurl,hyperref}
\setlength{\parindent}{0pt}
\setlength{\parskip}{5pt}
\setlength{\emergencystretch}{3em}
\newcommand{\N}{\mathbb{N}}
\newcommand{\Z}{\mathbb{Z}}
\newcommand{\Q}{\mathbb{Q}}
\newcommand{\R}{\mathbb{R}}
\newcommand{\C}{\mathbb{C}}
\newcommand{\F}{\mathbb{F}}
\newcommand{\A}{\mathbb{A}}
\newcommand{\PP}{\mathbb{P}}
\newcommand{\E}{\mathbb{E}}
\newcommand{\eps}{\varepsilon}
\newcommand{\dd}{\,\mathrm{d}}
\begin{document}
\section*{UnsolvedMath ID: 700737}
\textbf{Source problem code:} catalogue entry
\subsection{Definitions and mathematical statement}
Let \(X\) be a finite set of size \(v\), and write \(\binom{X}{k}\) for its set of \(k\)-element subsets. A \(6\)-\((v,k,\lambda)\) design is a pair \(\mathcal D=(X,\mathcal B)\), where \(\mathcal B\subseteq\binom{X}{k}\) is a set of blocks, without repetitions, and
\[
\left|\{B\in\mathcal B:T\subseteq B\}\right|=\lambda
\qquad\text{for every }T\in\binom{X}{6}.
\]
Its automorphism group consists of the permutations of \(X\) preserving \(\mathcal B\) setwise. A subgroup \(G\leq\operatorname{Aut}(\mathcal D)\) is block-transitive if, for every \(B_1,B_2\in\mathcal B\), some \(g\in G\) satisfies \(g(B_1)=B_2\). A flag is an incident pair \((x,B)\) with \(x\in B\in\mathcal B\); flag-transitivity means transitivity on these pairs.

\textbf{Unresolved qualification.} The supplied source does not define ``nontrivial.'' An explicit provisional interpretation is
\[
6<k<v,\qquad \mathcal B\neq\binom{X}{k}.
\]
The exclusion of complete designs and the boundary conditions in this interpretation require confirmation against the source's intended convention.

\textbf{Requested statement, under that provisional interpretation.} For every choice of integers \(v,k,\lambda\) with \(6<k<v\) and \(\lambda\geq1\), there is no \(6\)-\((v,k,\lambda)\) design \(\mathcal D=(X,\mathcal B)\) with \(\mathcal B\neq\binom{X}{k}\) admitting a block-transitive subgroup of \(\operatorname{Aut}(\mathcal D)\). This is the proposed nonexistence assertion, not a theorem established here.

\textbf{Supplied historical qualifications.} The original statement reports that the proposers had shown the corresponding nonexistence for nontrivial block-transitive \(8\)-designs, whereas some block-transitive, even flag-transitive, \(5\)-designs exist. These reports inherit the unresolved meaning of ``nontrivial.''

The comment labelled 2009 reports that a block-transitive group \(G\) on a nontrivial \(6\)-design must either be an affine group \(\operatorname{AGL}(d,2)\) or satisfy
\[
\operatorname{PSL}(2,q)\leq G\leq\operatorname{P\Gamma L}(2,q).
\]
It further attributes to M. Huber's 2010 paper the restriction that, when \(\lambda=1\), the only remaining possibility is \(G=\operatorname{P\Gamma L}(2,p^e)\), where \(p\in\{2,3\}\) and \(e\) is an odd prime power. Here an odd prime power means \(r^a\) for an odd prime \(r\) and an integer \(a\geq1\). The supplied excerpt does not specify further parameter or permutation-action qualifications.

The comment labelled 2013 reports that there are no block-transitive \(7\)-designs in the case \(\lambda=1\). Its sentence does not explicitly repeat a nontriviality hypothesis, so its intended scope requires checking.

The comment labelled 2021 reports that there are no nontrivial \(6\)-designs with \(k\leq10^4\) when the automorphism group is almost simple. A finite group \(H\) is almost simple if \(S\leq H\leq\operatorname{Aut}(S)\) for some nonabelian finite simple group \(S\). This comment does not explicitly repeat block-transitivity or specify whether it concerns the full automorphism group or a chosen block-transitive subgroup; no resolution of those scope issues is asserted here.
\subsection{Short English statement}
The problem asks whether every nontrivial finite \(6\)-design fails to admit a group transitive on its blocks. The source leaves ``nontrivial'' undefined; the explicit formulation above provisionally excludes \(k\leq6\), \(k=v\), and complete designs. The registry labels the problem open, and supplies several dated reports of partial results. Neither that status nor those reports have been independently verified.
\subsection{Sources}
\begin{itemize}
\item Supplied local UnsolvedMath registry record: permanent ID \(700737\), exact source code \texttt{catalogue entry}, title ``Kourovka Notebook Problem 11.45.'' The record attributes the problem to P. J. Cameron and C. E. Praeger, proposed in 1990, and labels it \texttt{open}. Its supplied update date is \texttt{2026-10-16}; this is registry metadata, not independent source verification.
\item \emph{The Kourovka Notebook (Unsolved Problems in Group Theory)}, original problem 11.45, 11th issue (1990), as supplied from the 21st-issue edition (2026), \texttt{arXiv:1401.0300v48}. The set description names E. I. Khukhro and V. D. Mazurov as editors. This attribution and edition information are supplied source data.
\item The comment labelled 2009 cites a contribution attributed in context to the proposers in \emph{Finite Geometry and Combinatorics (Deinze 1992)}, Cambridge University Press, 1993, pp. 103--119. The supplied record gives no contribution title. It attributes the affine/projective restriction on block-transitive groups to this reference.
\item M. Huber, \emph{Journal of Combinatorial Theory, Series A} \textbf{117}, no. 2 (2010), pp. 196--203, cited in the comment labelled 2009 for the \(\lambda=1\) restriction to \(\operatorname{P\Gamma L}(2,p^e)\). No article title is supplied.
\item M. Huber, \emph{Discrete Mathematics \& Theoretical Computer Science} \textbf{12}, no. 1 (2010), pp. 123--132, cited in the comment labelled 2013 for the reported nonexistence of block-transitive \(7\)-designs with \(\lambda=1\). No article title is supplied.
\item Q. Tan, W. Liu, J. Chen, \emph{Algebra Colloquium} \textbf{21}, no. 2 (2014), pp. 231--234, cited in the comment labelled 2021 for the reported \(k\leq10^4\) result involving an almost simple automorphism group. No article title is supplied.
\item All historical results and literature assessments above are reproduced as supplied attribution. The references have not been independently consulted or verified. No separate \texttt{literature\_checked\_at} date or \texttt{literature\_sources} links were supplied.
\end{itemize}
\end{document}
